\section{E. The next approximations and exactness on coverings}
\label{sec:E}

We write $F=\langle x,y\rangle$ for the discrete group, $F^{(i)}$
for its derived series, and $\gamma_iF$ for its lower central series,
with $F^{(0)}=\gamma_1F=F$. All commutators in this section are
$[a,b]=aba^{-1}b^{-1}$. The puncture used to define depth is the
puncture labelled $1$; it is additional data. The manuscript anchors
are \S51.16, folios 733v--736v, scans 172--178, especially (371)--(379).
The arithmetic model in those passages is
$X_n:x_0^n+x_1^n+x_\infty^n=0$, with $U_n$ obtained by deleting the
coordinate divisors. We retain that model over $\mathbb Q$.

\subsection{Derived quotients and the depth filtration}

\begin{statement}{E.1}{Proved}{The third derived quotient and its exact relation module}
Put $Q=F/F''$ and $R=\mathbb Z[Q]$. There is an exact sequence of
discrete groups and an exact sequence of left $R$-modules
\[
 1\longrightarrow F''/F'''\longrightarrow F/F'''
       \longrightarrow Q\longrightarrow1,
 \qquad
 0\longrightarrow F''/F'''\longrightarrow R^2
 \xrightarrow{d}R\xrightarrow{\epsilon}\mathbb Z\longrightarrow0,
\]
where $d(a,b)=a(\bar x-1)+b(\bar y-1)$, with right factors in this
formula and the regular left action on each term. This is an explicit
module description over the noncommutative ring $R$.
\end{statement}
\begin{proof}
The cover of the two-loop graph corresponding to the normal subgroup
$F''$ has vertices $Q$ and oriented edges $(q,x),(q,y)$ from $q$ to
$q\bar x,q\bar y$. Its cellular differential is the displayed map.
The graph is connected, so its reduced zeroth homology vanishes.
Its fundamental group is $F''$, and its first homology is
$F''/[F'',F'']=F''/F'''$. Lifting paths identifies conjugation by
$F$ with the left deck action of $Q$. These identifications give every
arrow, its kernel and its image. The group sequence is the quotient
sequence for the nested normal subgroups $F'''\subset F''$.

For the algebraic constructions below we use the rational relative
unipotent completion over a specified finite quotient, constructed in
E.14. For an unbased finite-dimensional representation problem one
can also use the affine proalgebraic group with coordinate Hopf algebra
the matrix coefficients of all finite-dimensional rational
representations of $F/F'''$. Evaluation is the canonical group map
into its rational points, and factoring matrix coefficients proves
its universal property. These constructions specify an envelope and
its comparison maps; they do not require a finite-index abelian
subgroup of $F/F'''$.
\end{proof}

\begin{statement}{E.2}{Proved}{The depth-two Lie algebra and its comparison with the derived quotient}
Let $L=\mathbb L_{\mathbb Q}(x,y)$, graded by total degree, and let
$D^rL$ be the span of Lie monomials with at least $r$ occurrences of
$y$. Put $v_i=\operatorname{ad}_x^i(y)$ and
$V=\bigoplus_{i\geq0}\mathbb Qv_i$. There is an isomorphism of graded
Lie algebras
\[
 L/D^3L\simeq \mathbb Qx\ltimes(V\oplus\Lambda^2V).
\]
Here $[v_i,v_j]=v_i\wedge v_j$, the bracket of $V$ with
$\Lambda^2V$ vanishes, and $[x,v_i]=v_{i+1}$; the action of $x$
on $\Lambda^2V$ is the induced derivation. In weight $n\geq2$ the
depth-two dimension is $\lfloor(n-1)/2\rfloor$.
The identity on $x,y$ induces the surjection
$L/L'''\twoheadrightarrow L/D^3L$ with kernel $D^3L/L'''$.
\end{statement}
\begin{proof}
We give the elimination argument. Jacobi moves every occurrence of
$x$ in a bracket in the ideal generated by $y$ onto a bracket entry.
Consequently that ideal is generated as a Lie algebra by the $v_i$.
In the tensor algebra, the leading words of the $v_i$ are $x^iy$.
The words obtained by concatenating these blocks have unique
factorizations at the occurrences of $y$. Thus the homomorphism
from the tensor algebra on the symbols $v_i$ into
$\mathbb Q\langle x,y\rangle$ is injective, by its triangular leading
words. Its primitive Lie subalgebra is the free Lie algebra on the
$v_i$. Keeping at most two such symbols yields $V\oplus\Lambda^2V$.
The derivation formula follows from Jacobi. A basis in weight $n$
is $v_i\wedge v_j$ with $i<j$ and $i+j=n-2$, which has the stated
cardinality.

Every element of $L'$ has depth at least one, hence
$L''\subset D^2L$ and $L'''\subset D^4L\subset D^3L$.
This proves the displayed quotient map and its kernel.
The element $[y,[y,[y,x]]]$ has weight four and belongs to $D^3L$;
its tensor expansion is nonzero, whereas $L'''$ has no weight-four
part. Thus the kernel is nonzero. Also
$[[x,y],[x,[x,y]]]$ survives in $\Lambda^2V$, so the identity on
$x,y$ does not descend from $L/L''$ to $L/D^3L$: its defining
kernel $L''$ contains this explicitly nonzero image.
Completion for the total-degree filtration gives the same maps of
pronilpotent Lie algebras, and exponentiation gives maps of rational
prounipotent groups. No integral exponential is used.
\end{proof}

\begin{statement}{E.3}{Known}{The rational Galois calculation in depth two through weight fourteen}
Normalize the depth-one generators by
$\bar\sigma_{2m+1}=\operatorname{ad}_x^{2m}(y)$.
For the positive-weight Lie algebra of the outer pro-$\ell$ Galois
image, tensored with $\mathbb Q_\ell$, the depth-one dimensions are
one in weights $3,5,7,9,11,13$ and zero in the other weights at most
fourteen. Its depth-two dimensions in weights $1,\ldots,14$ are
\[
 0,0,0,0,0,0,0,1,0,1,0,1,0,2.
\]
The first relation among the depth-one brackets is
\[
 \{\bar\sigma_3,\bar\sigma_9\}
       -3\{\bar\sigma_5,\bar\sigma_7\}=0
 \quad\text{in depth two}.
\]
This is a statement after tensoring with $\mathbb Q_\ell$ and taking
the specified filtrations; it asserts no integral saturation index.
\end{statement}
\begin{proof}
We use Hain--Matsumoto, Theorem 8.4, its application in \S8.4,
Lemma 8.5 and Remark 8.6 of the opened version
\cite{E_HainMatsumoto}, and Brown's depth-two bracket formula and
exact sequence (7.8), Example 7.2 \cite{E_BrownDepth}.
Here $U_1$ has $H^1_{\mathrm{et}}(U_{1,\overline{\mathbb Q}},
\mathbb Q_\ell(1))=\mathbb Q_\ell^2$ with trivial Galois action:
the two Kummer covers are those of $t$ and $1-t$.
Its outer pro-$\ell$ action is unramified away from $\ell$.
The indicated theorem therefore identifies the weight-graded Lie
algebra of its Zariski closure with the rationalized graded Galois
image. The weight-one generator maps to zero by Lemma 8.5;
the odd generators in weights at least three have nonzero
depth-one coefficient, as recorded in Remark 8.6. This supplies
actual Galois elements after rationalization, rather than only
solutions of a linear equation.

The polynomial form of the depth-two bracket is
\begin{align*}
 B_{a,b}(X,Y)={}&X^a\bigl(Y^b-(Y-X)^b\bigr)
 +(Y-X)^a\bigl(X^b-Y^b\bigr)\\
 &+Y^a\bigl((Y-X)^b-X^b\bigr),
 \qquad a,b\text{ positive and even}.
\end{align*}
One can derive this formula directly by substituting
$D_f(x)=0$, $D_f(y)=[y,f]$ into
$[f,g]+D_f(g)-D_g(f)$ in the tensor algebra, and reading the
two $y$-blocks from the right. The program \texttt{e\_depth.py}
does both expansions independently over the integers and compares
every coefficient. The only pairs in weights $8,10,12,14$ are,
respectively, $(3,5)$, $(3,7)$, $(3,9),(5,7)$, and
$(3,11),(5,9)$. Their exact ranks are $1,1,1,2$.
The other weights contain no pair.

For completeness of this span in depth two, every bracket word of
length at least three in the Galois generators has depth at least
three. A generator itself has a nonzero depth-one image in its
weight; different generator weights cannot cancel in a weight-graded
piece. Thus the generation theorem reduces the calculation exactly
to the displayed pairs. The polynomial
\[
 X^2Y^2(X-Y)^3(X+Y)^3
 =X^8Y^2-3X^6Y^4+3X^4Y^6-X^2Y^8
\]
satisfies the antisymmetry and the three-term period equation.
Substitution in $B_{a,b}$ gives the displayed relation, also checked
word by word in the free tensor algebra. Brown's exact sequence
identifies its line with the weight-twelve cusp-form period line.
This verifies the requested modular relation without making a
freeness assertion for an integral Galois Lie algebra.
\end{proof}

\begin{statement}{E.4}{Proved}{What the metabelian representation forgets about these brackets}
Under the canonical Lie homomorphism $L\to L/L''$, every Ihara
bracket of two elements of $L'$ maps to zero. The depth-two image
of $\{\bar\sigma_3,\bar\sigma_5\}$ is nonzero. Thus the specified
comparison maps exhibit information already added in weight eight;
the first dependence between these new bracket directions occurs
in weight twelve as in E.3.
\end{statement}
\begin{proof}
Write the derived module of $L/L''$ as
$\mathbb Q[X,Y]c$, where $c=[x,y]$ and $X,Y$ act by
$\operatorname{ad}_x,\operatorname{ad}_y$. For $f=pc$ one has
$D_f(c)=[x,[y,pc]]=XYpc$. The derivation $D_f$ vanishes on the
operators $X,Y$ when applied to the derived module: its correction
there is a bracket of two derived elements. Hence
$D_f(qc)=XYpq\,c=D_g(pc)$ for $g=qc$, while $[f,g]=0$.
The Ihara bracket therefore vanishes. In E.3 the polynomial
$B_{2,4}=2X^5Y-5X^4Y^2+5X^2Y^4-2XY^5$ is explicitly nonzero.
Both claims concern the images of the same derivations through the
two quotient maps just constructed.
\end{proof}

\subsection{Localizers and finite coverings}

\begin{statement}{E.5}{Proved}{Two-out-of-three and relative Theorem A}
For each of abelianization, the metabelian quotient, a fixed
nilpotent quotient $G/\gamma_{c+1}G$ with $c\geq1$, pro-$p$
completion, and rational prounipotent completion, the class
$W_{\mathfrak a}$ in the question satisfies two-out-of-three,
retract closure, the terminal-object axiom and relative Theorem A
for all small categories.
\end{statement}
\begin{proof}
We first specify the tests determining each class. They are the
classifying groupoids $BH$ for, respectively, all abelian groups,
all metabelian groups, all groups of nilpotence class at most $c$,
all finite $p$-groups, and all groups $U(\mathbb Q)$ of rational
points of finite-dimensional unipotent rational algebraic groups.
We include all discrete sets as tests to detect components.
For a category $C$, functors $C\to BH$ invert every arrow and
therefore factor uniquely through its fundamental groupoid.
For a connected pointed category the homotopy fiber of evaluation
$\operatorname{Fun}(C,BH)\to BH$ is equivalent to the discrete groupoid
$\operatorname{Hom}(\pi_1C,H)$. The universal properties of the five
approximations, followed by Yoneda, prove that a component-bijective
functor is in $W_{\mathfrak a}$ precisely when it induces an
equivalence of these functor groupoids for every indicated test.
Conjugations cause no loss here: the evaluation fiber restores the
pointed homomorphisms. Two-out-of-three and retract closure follow
from equivalences of groupoids. A category with a terminal object
has contractible nerve and passes every test.

Here is the relative argument. For $C\to S$, the Grothendieck
construction of $s\mapsto C/s$ has objects
$(s,c,u:f(c)\to s)$. Its forgetful functor to $C$ has a section
$c\mapsto(f(c),c,1)$ and a natural transformation from section
followed by forgetful to the identity, furnished by $u$.
It therefore has the same fundamental groupoid as $C$.
Functors from this construction to a groupoid $T$ consist exactly
of functors $C/s\to T$ with invertible coherent transition
transformations along the arrows of $S$. This is the pseudolimit
of the groupoids $\operatorname{Fun}(C/s,T)$. A pointwise
equivalence of such diagrams induces an equivalence of pseudolimits:
choose objectwise lifts, transport the transitions through full
faithfulness, and use faithfulness to verify their coherence;
the same argument applies to morphisms. A map $C\to D$ over $S$
which passes every test on every comma category consequently passes
the test on $C\to D$. This is the stated relative Theorem A.
\end{proof}

\begin{statement}{E.6}{Proved}{The exact comparison maps for quotient approximations}
Let $q:F\twoheadrightarrow G$ be any finite quotient, $N=\ker q$,
and $V(H)$ one of $H'$, $H''$, or $\gamma_{c+1}H$.
There is a natural short exact sequence
\[
 1\longrightarrow\frac{N\cap V(F)}{V(N)}
 \longrightarrow\frac{N}{V(N)}
 \longrightarrow
 \ker\left(\frac{F}{V(F)}\longrightarrow\frac{G}{V(G)}\right)
 \longrightarrow1.
\]
The comparison requested with target $G$ itself is the specialization
$V(G)=1$ of this exact sequence. The map $F/V(F)\to G$ with
composite $q$ exists exactly when $V(G)=1$.
\end{statement}
\begin{proof}
Surjectivity of $q$ gives $q(V(F))=V(G)$, by the definitions of
the verbal subgroups. The inverse image of $V(G)$ is $NV(F)$,
so the kernel in the display is $NV(F)/V(F)$. The map from $N$
onto it has kernel $N\cap V(F)$ and kills $V(N)$.
This proves exactness and identifies every map. A factorization
of $q$ through $F/V(F)$ kills $q(V(F))=V(G)$; the converse is
the same calculation. The formulas thus also give an explicit
replacement target for every quotient for which the arrow to $G$
written in the question does not descend.
\end{proof}

\begin{statement}{E.7}{Proved}{All nontrivial finite quotients have a nonzero comparison kernel}
For every finite quotient $F\twoheadrightarrow G$ of order $d>1$,
each of the three comparison kernels in E.6 is nonzero.
For abelianization it is a free abelian group of rank $d-1$;
the image is a rank-two lattice of index $|G^{\mathrm{ab}}|$ in
$F^{\mathrm{ab}}$. For the nilpotent quotient of class $c$ its
rational Hirsch rank is
\[
 \sum_{j=1}^c\bigl(w_j(d+1)-w_j(2)\bigr),\qquad
 w_j(r)=\frac1j\sum_{e\mid j}\mu(e)r^{j/e}.
\]
\end{statement}
\begin{proof}
The connected Cayley graph of $G$ has $d$ vertices and $2d$ edges.
A maximal tree gives $d+1$ free generators of $N$, hence
$N^{\mathrm{ab}}=\mathbb Z^{d+1}$. The abelian image and its index
follow from E.6. A subgroup of a free abelian group is free, and
the rank of the kernel is $(d+1)-2=d-1$.

For the nilpotent case the domain is the free nilpotent group on
$d+1$ generators and the image is a finite-index subgroup of the
free nilpotent group on two generators. Both are torsion-free.
Their lower central ranks are the Witt numbers: the tensor algebra
identity $1/(1-rt)=\prod_{j\geq1}(1-t^j)^{-w_j(r)}$, obtained by
ordering a homogeneous Lie basis in the universal enveloping
algebra, gives the displayed formula by logarithms and M\"obius
inversion. Hirsch rank is additive in a short exact sequence of
torsion-free nilpotent groups. The first summand is already $d-1$.

For the metabelian case choose
$a\in N\cap F'$ with nonzero image $A$ in $N^{\mathrm{ab}}$;
such an element exists by the abelian calculation. Choose $t\in N$
whose abelian class $T$ is linearly independent from $A$ over
$\mathbb Q$. In the free metabelian quotient of $N$ the module
$N'/N''$ injects into $\mathbb Z[N^{\mathrm{ab}}]^{d+1}$ by the
covering-graph differential. It is therefore torsion-free over this
integral domain. The image of $[t,a]$ is nonzero: its degree-two
term is $T\wedge A$. The element
\[
 w=[a,tat^{-1}]\in N\cap F''
\]
equals $[a,[t,a]]$ and has image $(A-1)[t,a]$ in that module
for our convention $[u,v]=uvu^{-1}v^{-1}$. This image is nonzero.
Thus $w\notin N''$, giving a specific construction of a nontrivial
element of $(N\cap F'')/N''$ for every such quotient.
For $d=1$, $N=F$ and every comparison is the identity.
\end{proof}

\begin{statement}{E.8}{Proved}{The pro-$p$ comparison and the precise exactness criterion}
Write $G^{(p)}$ for the maximal $p$-group quotient of the finite
group $G$. There is a continuous surjection
\[
 \widehat N_p\longrightarrow
 \ker(\widehat F_p\longrightarrow G^{(p)}).
\]
Its source is free pro-$p$ of rank $|G|+1$ and its target is free
pro-$p$ of rank $|G^{(p)}|+1$. It is an isomorphism exactly for
$p$-groups $G$. For $p$-groups the target is the open subgroup
corresponding to the given cover.
More explicitly, put $r=|G|+1$ and $s=|G^{(p)}|+1$.
There are free pro-$p$ bases $a_1,\ldots,a_s,b_1,\ldots,b_{r-s}$
in the source and $\bar a_1,\ldots,\bar a_s$ in the target
for which the comparison sends $a_i\mapsto\bar a_i$ and
$b_j\mapsto1$. Its kernel is exactly the closed normal subgroup
generated by those $r-s$ elements $b_j$. This description is
noncanonical; the original continuous comparison is canonical.
\end{statement}
\begin{proof}
The image is the closure of $N$. The quotient by that closure is
the maximal pro-$p$ quotient of the finite group $F/N$, namely
$G^{(p)}$. This proves surjectivity. For a normal subgroup of
$p$-power index the induced pro-$p$ topology is its full topology:
the core in the ambient group of a normal subgroup of finite
$p$-power index is the intersection of finitely many conjugates,
and its quotient is an extension of finite $p$-groups.
Completing therefore identifies $\widehat N_p$ with the open
subgroup when $G$ is a $p$-group. The Schreier rank calculation
applied to finite $p$-quotients, or the covering tree, gives the
two ranks. An isomorphism would identify their Frattini quotients,
whose dimensions are these ranks, and hence forces
$|G|=|G^{(p)}|$. This equality says precisely that $G$ is a
$p$-group. All arrows used here are continuous homomorphisms of
compact pro-$p$ groups.

For the kernel description, lift a target basis to elements
$a_1,\ldots,a_s$. Their classes in the source Frattini quotient
are independent. Extend them to a basis of that vector space and
choose lifts $b_j$. The pro-$p$ basis theorem and Hopficity show
that the resulting $r$ elements form a free basis. The subgroup
generated by the $a_i$ maps isomorphically onto the target:
the induced surjection of free groups of rank $s$ is an
isomorphism by Hopficity. For each $b_j$, choose the unique
word $w_j(a_1,\ldots,a_s)$ with its target image and replace
$b_j$ by $b_jw_j^{-1}$. The Frattini classes remain a basis,
so these are again free generators. The quotient by their closed
normal closure is now the free group on the $a_i$, proving the
stated kernel and image, not merely their ranks.
\end{proof}

\begin{statement}{E.9}{Proved}{Ordinary rational completion of a finite-index subgroup}
For $|G|=d$, the canonical morphism of rational prounipotent groups
$N^{\mathrm{un}}\to F^{\mathrm{un}}$ is surjective. Its map on
abelianized Lie algebras is a surjection
$\mathbb Q^{d+1}\to\mathbb Q^2$ with kernel of dimension $d-1$.
It is an isomorphism exactly for $d=1$.
In suitable noncanonical free pronilpotent coordinates, its Lie
algebra map sends $X,Y$ to the two target generators and
$z_1,\ldots,z_{d-1}$ to zero. Its kernel is the closed ideal
generated by the $z_j$; by free-Lie elimination it is the completed
free Lie algebra on
$\mathbb Q\langle X,Y\rangle\otimes
\langle z_1,\ldots,z_{d-1}\rangle$, where a word acts by iterated
adjoints. Exponentiation is the exact kernel group.
Every morphism from a connected unipotent rational algebraic group
to the constant finite group $G$ is trivial; thus the specified
map to $G$ in a kernel formula descends only for the trivial
quotient. The preceding surjection is the comparison obtained
from the trivial completion of $G$.
\end{statement}
\begin{proof}
A homomorphism from $F$ to a unipotent group has the same Zariski
closure on $N$ as on $F$: the quotient of the closures has finite
order, while a unipotent group in characteristic zero is connected
and admits no nontrivial finite quotient. Equivalently, every
$f\in F$ has a positive power in $N$, and taking roots in the
one-parameter unipotent subgroup containing its image recovers
that image. The resulting map of completions is surjective.
Both Lie algebras are completed free Lie algebras on their
abelianizations, of ranks $d+1$ and two. The abelian comparison
is the rationalization of E.7. A morphism from a connected scheme
to a constant finite scheme is constant, which proves the last
assertion and types the replacement map.

For the coordinate assertion, lift the target free Lie generators
and complete their independent linear terms to a basis of the
source abelianization. The homomorphism of completed free Lie
algebras determined by these lifts is invertible: its map in
each graded degree is invertible, and degreewise inversion
converges. Its first two generators give a section of the target
free Lie algebra. Subtract from every remaining generator the
section of its image. This triangular change is again invertible
and puts that image equal to zero. The kernel is consequently
the closed ideal generated by the remaining $z_j$.
Jacobi writes this ideal as the Lie algebra generated by the
iterated $X,Y$ adjoints of the $z_j$. In the free associative
algebra their leading words have unique cuts at the letters
$z_j$, proving freeness exactly as in E.2. This also computes
the whole kernel for every finite quotient, including those for
which the proposed map to the original $G$ has no descent.
\end{proof}

\begin{statement}{E.10}{Proved}{Which finite regular coverings preserve the localizers}
For a regular covering with monodromy $q:F\twoheadrightarrow G$,
membership in $\operatorname{Fib}_{W_{\mathfrak a}}$ holds only for
$G=1$ for abelianization, the metabelian quotient, every fixed
positive nilpotent class, and rational prounipotent completion.
For pro-$p$ completion it holds exactly for finite $p$-groups $G$.
These are exactly the cases of isomorphism in E.7--E.9, with the
properly descended target of E.6 or E.8.
\end{statement}
\begin{proof}
We first give uniform counterexamples over $BF$, so the universal
base-change requirement is tested, not merely the given fiber.
For a nontrivial finite $G$, choose a prime $r\nmid |G|$ and let
$V$ be the augmentation-zero summand of the regular
$\mathbb F_r[G]$-module. Averaging shows $V_G=0$, and $V\ne0$.
Let $A=V\rtimes F$ for the action through $q$, with retraction
$A\to F$. The inclusion $F\to A$ is an abelian and every
nilpotent equivalence: $[F,V]=V$ puts $V$ in every term
$\gamma_jA$, $j\geq2$. It is also a rational unipotent
equivalence, since the finite normal subgroup $V$ maps trivially
to every rational unipotent group. Pullback along the $G$-cover
replaces the inclusion by $N\to V\times N$. Its abelian and
nilpotent quotients retain $V$, so they fail those tests.
For the rational test use instead the integral augmentation ideal
$V=\ker(\mathbb Z[G]\to\mathbb Z)$. Every homomorphism
$A=V\rtimes F\to U(\mathbb Q)$, with $U$ a rational unipotent
group, kills $V$. To prove this, take the Zariski closure of the
image of $V$. It is an abelian unipotent group. Conjugation by
each image of an element of $F$ acts unipotently on its Lie
algebra and has finite order, since the action on $V$ factors
through $G$. A finite-order unipotent linear operator in
characteristic zero is the identity. Hence the image factors
through $V_G$, which is finite because its rationalization is
zero. A unipotent group in characteristic zero has no nontrivial
finite subgroup, so that image is trivial. The universal property
therefore identifies the rational prounipotent completions of
$F$ and $A$. After pullback, $N\to V\times N$ retains the
nonzero vector quotient $V\otimes\mathbb Q$, of dimension $d-1$.

For the metabelian test form a finite quotient
$H=F/(N'N^r)$, with $r$ any prime. The inclusion
$N\cap F'\hookrightarrow N$ has nonzero primitive abelian
image of rank $d-1$, by E.7. Its reduction modulo $r$ remains
nonzero. Thus $H'$ is nontrivial and the image of $N'$ is trivial.
Choose a prime $s\nmid |H|$ and put
$V=(1-e_{H'})\mathbb F_s[H]$, where
$e_{H'}=|H'|^{-1}\sum_{h\in H'}h$. This is a nonzero $H$-module
with $[H',V]=V$. For $A=V\rtimes F$ one has
$A''=V\rtimes F''$, so $F\to A$ induces an isomorphism of
metabelian quotients. On the pulled-back group $V\rtimes N$,
$N'$ acts trivially on $V$; its second derived subgroup is
$\{0\}\rtimes N''$. Consequently the metabelian quotient
retains $V$, and $N\to V\rtimes N$ fails the test. These are
explicit finitely presented group constructions. Their standard
finite presentation complexes and mapping cylinders give finite
simplicial-category models with the same fundamental-group
calculation.

For a finite $p$-group cover, E.8 applies to every group occurring
after any base change, not just $F$. An isomorphism of pro-$p$
completions identifies the finite orbit sets and the appropriate
open stabilizer subgroups, whose completions are those open
subgroups. Hence pullback preserves $W_p$.

For a regular cover with $G$ not a $p$-group, let $P$ be a Sylow
$p$-subgroup. The image of
$F\to G\times(F/F'F^p)$ is the fiber product over
$G/G'G^p$: this follows by taking the two projections and their
kernels, or directly from $q(F'F^p)=G'G^p$.
The map $P\to G/G'G^p$ is surjective, because its cokernel has
order both a power of $p$ and prime to $p$. Choose words $u,v$
mapping into $P$ whose images in $F/F'F^p$ are $x,y$.
The endomorphism $x\mapsto u,y\mapsto v$ induces a pro-$p$
automorphism: it is onto the Frattini quotient, hence onto the
finitely generated free pro-$p$ group, and that group is Hopfian.
Its pulled-back regular $G$-set has image contained in the proper
subgroup $P$, and is therefore disconnected. This violates the
component condition in $W_p$. Finally, a trivial covering is a
product with a discrete set and preserves all five classes.
\end{proof}

\begin{statement}{E.11}{Proved}{The explicit three-letter obstruction in every nilpotent approximation}
The endomorphism $\alpha(x)=x$, $\alpha(y)=y[x,y]$ induces the
identity on every graded lower central quotient and an
automorphism on every nilpotent quotient, every pro-$p$
completion and the rational prounipotent completion of $F$.
Pullback of the action $x\mapsto(12),y\mapsto(13)$ changes its
orbit set from one orbit to $\{1,2\},\{3\}$.
It consequently refutes intrinsic-fibration membership for this
cover for each of these approximations, and for every functor
of their inverse system which sends isomorphisms to isomorphisms.
\end{statement}
\begin{proof}
Substituting a generator times a commutator in an iterated
commutator changes it only in the next lower central degree.
Thus the graded action is the identity. Induction in the central
extensions gives an automorphism on each nilpotent quotient.
The Frattini argument of E.10 proves the pro-$p$ statement;
the completed free rational Lie algebra gives the unipotent
statement by degreewise inversion. The inverses commute with
the transition maps, so the map of inverse systems is an
isomorphism. For $a=(12),b=(13)$, multiplication gives
$b[a,b]=a$. Both pulled-back generators therefore act by
$(12)$. The program \texttt{e\_obstructions.py} verifies these
permutations and the orbit sets. The ordinary metabelian quotient
retains the map to $S_3$, since $S_3''=1$; the map of E.11 is
not substituted for a metabelian equivalence.
\end{proof}

\begin{statement}{E.12}{Proved}{Exhaustive small-quotient and Schreier verification}
For all ordered generating pairs in every group model of order
at most twenty-four, the program \texttt{e\_smallgroups.py}
constructs an integral cycle basis of $N^{\mathrm{ab}}$, the
comparison to $F^{\mathrm{ab}}$, and the deck character; it
computes the free-metabelian relation-module generic rank and
the first four lower central ranks. Every generated kernel has
rank $|G|+1$, the comparison lattice has index
$|G^{\mathrm{ab}}|$, and the rational deck character is
$1+\operatorname{reg}_G$. All checks pass.
\end{statement}
\begin{proof}
The input is the complete order-at-most-twenty-four list in the
opened primary GroupNames database \cite{E_GroupNames}.
It contains seventy-four groups. The accompanying input file
records the permutation generators and exact source URLs.
The program closes each set of generators, checks its order,
constructs its multiplication table, and enumerates every ordered
pair. It tests generated subgroup size without a presentation
heuristic. For every generating pair it constructs the $2|G|$
edges and a spanning tree, and checks that every fundamental
cycle has zero boundary. The non-tree edge coefficients form
an integral basis, so this proves the integral calculation.
The gcd of all two-by-two minors of the pushforward matrix
checks the Smith index. Deck traces are calculated for every
group element. Independent tests exhaust the images of a
generating pair in every other two-generated model of the same
order to verify pairwise nonisomorphism of the models used.

The lower central ranks are calculated from the proved Witt
formula in E.7. For the free group $N$ on $r=|G|+1$ generators,
the graph over $N^{\mathrm{ab}}$ gives
$N'/N''=\ker(R_N^r\to R_N)$, $R_N=\mathbb Z[N^{\mathrm{ab}}]$;
localization to its fraction field has dimension $r-1=|G|$.
Thus the printed free-metabelian invariant has an exact module
presentation behind it. The output is accompanied by
\texttt{e\_all\_surjections.csv}, with a row for every ordered
epimorphism, and integral cycle witnesses for every group.
The formulas E.6--E.10 prove the assertions for all finite
quotients, so the bounded computation also has a general proof.
\end{proof}

\begin{statement}{E.13}{Proved}{Specialization to the Fermat coverings}
For $N_n=\ker(F\to(\mathbb Z/n)^2)$ the abelian comparison is
\[
 0\longrightarrow\mathbb Z^{n^2-1}
 \longrightarrow N_n^{\mathrm{ab}}
 \longrightarrow n\mathbb Z^2\longrightarrow0.
\]
The metabelian and class-$c$ comparisons are the exact sequences
of E.6 with kernels $F''/N_n''$ and
$\gamma_{c+1}F/\gamma_{c+1}N_n$, respectively; both are nonzero
for $n>1$. Pro-$p$ completion is exact exactly when $n$ is a
power of $p$. With $n=p^a m$, $(m,p)=1$, its replacement target
is the open subgroup of $\widehat F_p$ of index $p^{2a}$, and
the source has rank $n^2+1$. The rational unipotent comparison
is the surjection of E.9, with abelian kernel dimension $n^2-1$.
\end{statement}
\begin{proof}
The abelian covering map sends a closed edge cycle to its two
exponent sums. Each sum is divisible by $n$, and the loops
$x^n,y^n$ realize the two generators of $n\mathbb Z^2$.
The source rank is $n^2+1$ by the graph calculation. Since
$N_n$ contains $F'$, it contains both $F''$ and every
$\gamma_{c+1}F$ for $c\geq1$, giving exactly the two kernels
in the statement. E.7 proves their nonvanishing. The maximal
$p$-quotient of $(\mathbb Z/n)^2$ is $(\mathbb Z/p^a)^2$,
so E.8 gives the asserted ranks and the exactness criterion.
E.9 gives the last assertion. These are the maps induced by
the actual covering inclusion, and coincide with the Fox
matrices computed in section B.
\end{proof}

\begin{statement}{E.14}{Proved}{Relative completion retains the entire covering kernel}
The rational relative unipotent completion $\mathcal R_q$ of
$F\twoheadrightarrow G$, for finite $G$, fits into
\[
 1\longrightarrow N^{\mathrm{un}}
 \longrightarrow\mathcal R_q\longrightarrow G_{\mathbb Q}
 \longrightarrow1.
\]
It splits over $\mathbb Q$, with splittings conjugate by
$N^{\mathrm{un}}(\mathbb Q)$. Its quotient with abelian kernel
has vector kernel $H_1(N,\mathbb Q)$, of dimension $|G|+1$ and
character $1+\operatorname{reg}_G$.
For $q_2:F\to(\mathbb Z/2)^2$ its four character
multiplicities are $2,1,1,1$.
\end{statement}
\begin{proof}
Choose a transversal of $G$ in $F$. Its conjugation maps on $N$
extend uniquely to $N^{\mathrm{un}}$, and its products give a
factor set with values in $N\to N^{\mathrm{un}}$.
The associated multiplication on the disjoint union of $|G|$
copies of $N^{\mathrm{un}}$ is associative by the factor-set
identities. This constructs $\mathcal R_q$ and the displayed
exact sequence as rational proalgebraic group schemes.
Every map from $F$ into an extension of $G$ by a unipotent group
extends uniquely on $N$ and then on the chosen components.
It respects multiplication by the same identities. This proves
the relative universal property.

At each successive central vector quotient of the unipotent
kernel, a section lifts because the cohomology in positive
degree of the finite group $G$ with rational coefficients
vanishes. To see the vanishing, average a cochain contraction
over $G$ and divide by $|G|$ in the standard bar resolution.
The same averaging in degree one conjugates two sections.
Induction gives compatible sections and conjugators through
the central quotients, and their inverse limits give the
assertions for $N^{\mathrm{un}}$. This specifies the conjugacy
ambiguity; uniqueness of the conjugating element is not claimed.
Finally, the chain modules of the finite Cayley graph are
$\mathbb Q[G]^2$ and $\mathbb Q[G]$, with zeroth homology
$\mathbb Q$. Semisimplicity of $\mathbb Q[G]$ gives
$H_1(N,\mathbb Q)\simeq\mathbb Q\oplus\mathbb Q[G]$.
For $G=(\mathbb Z/2)^2$, expanding its regular character gives
the four stated multiplicities. This action has a $-1$
eigenvalue; it survives in the relative completion, whereas
the map of E.9 takes coinvariants and loses these characters.
\end{proof}

\begin{statement}{E.15}{Proved}{The genuine action on a relative completion}
Let $\operatorname{Aut}(F,q)$ consist of the automorphisms of $F$
carrying $N$ to itself. Its homomorphism to
$\operatorname{Aut}_{\mathbb Q}(\mathcal R_q)$ is injective,
and its image is exactly the subgroup preserving the embedded
discrete subgroup $F\subset\mathcal R_q(\mathbb Q)$ setwise.
The automorphism $\operatorname{Int}([x,y])$, at $q=q_2$, acts
nontrivially on the relative completion, although its action
on $F^{\mathrm{ab}}$ and the image of its conjugator there are
both the identity.
\end{statement}
\begin{proof}
The universal property in E.14 extends an automorphism and its
inverse, with their induced automorphism of $G$, so the
extension is an automorphism. The map $N\to N^{\mathrm{un}}$
is injective. Write a nontrivial reduced word in maximal
nonzero-power blocks $z_{i_1}^{a_1}\cdots z_{i_k}^{a_k}$,
with adjacent $i_j$ different. Under $z_i\mapsto1+Z_i$,
the coefficient of $Z_{i_1}\cdots Z_{i_k}$ with these
$k$ alternating blocks is $a_1\cdots a_k\ne0$:
omitting a block cannot yield that sequence of $k$
blocks, and each retained block must contribute its
linear term. A finite truncation of the power-series
ring detects the word in its unipotent group $1+(Z_i)$.
Since
components detect $F/N$, the map $F\to\mathcal R_q(\mathbb Q)$
is injective as well. This proves faithfulness. An algebraic
automorphism preserving this subgroup restricts to an
automorphism of $F$ preserving its identity component $N$;
the universal property identifies its extension with the
given automorphism. This proves the image description.

At $q_2$ the commutator $w=[x,y]$ has a nonzero cycle class in
$H_1(N_2,\mathbb Q)$, as can be read from the unit-square
cycle in the four-vertex graph. Choose another cycle with
independent class. Their bracket is nonzero in the degree-two
free Lie algebra of $N_2^{\mathrm{un}}$, so conjugation by $w$
is nontrivial. Passing from a genuine action to its outer
class kills this particular inner action. Thus this example
also specifies the actual quotient map responsible for the
ambiguity, rather than identifying an outer action with a
genuine one.
\end{proof}

\subsection{Sections and arithmetic obstructions}

\begin{statement}{E.16}{Proved}{All sections of the rank-two toric quotient}
Let $T=\mathbb G_m^{\{0,1,\infty\}}/\mathbb G_m$ over
$\mathbb Q$, with its rational identity. Its arithmetic
fundamental-group extension has geometric kernel
$\widehat{\mathbb Z}(1)^2$. Its continuous sections modulo
conjugacy by this kernel form the abelian group
\[
 \left(\varprojlim_m\mathbb Q^*/\mathbb Q^{*m}\right)^2.
\]
One factor has the explicit description
\[
 \{\pm1\}\times
 \left\{(a_p)_p\in\prod_{p\text{ prime}}\widehat{\mathbb Z}:
 \forall m,\ a_p\in m\widehat{\mathbb Z}
 \text{ for all but finitely many }p\right\}.
\]
This computes the cohomology in (379) for its specified
rank-two quotient at every level $n$.
\end{statement}
\begin{proof}
The section at the identity gives the semidirect product
$\widehat{\mathbb Z}(1)^2\rtimes G_{\mathbb Q}$.
A section has form $g\mapsto(c(g),g)$; multiplication says
$c(gh)=c(g)+g c(h)$. Conjugation changes $c$ by a
coboundary. At finite level the Kummer sequence and
Hilbert 90 identify $H^1(G_{\mathbb Q},\mu_m)$ with
$\mathbb Q^*/\mathbb Q^{*m}$. To verify the reduction,
a multiplicative cocycle of a finite Galois extension
is a coboundary by choosing $b$ for which
$\sum_\tau c(\tau)\tau(b)\ne0$ and translating the sum;
linear independence of field embeddings supplies such
a $b$. Every algebraic cocycle is defined over a finite
extension. Taking limits gives the formula. The possible
$\varprojlim^1 H^0(\mu_m)$ vanishes because these groups
are finite and their descending images stabilize.

Unique prime factorization gives
$\mathbb Q^*=\{\pm1\}\times\bigoplus_p p^{\mathbb Z}$.
An element modulo $m$ has finitely many nonzero valuation
residues. Compatible residues are exactly the families
in the second display. Even $m$ retain the sign and odd
$m$ kill it; their compatible limit is $\{\pm1\}$.

The map $U_n\to T$ is inclusion on the open with all
coordinates nonzero. On geometric groups it is the
covering pushforward of E.13 divided by $n$, hence
surjective onto $\widehat{\mathbb Z}(1)^2$. The map to
the base torus in coordinates is
$(u,v)\mapsto(-u^n,-v^n)$, with
$t=-u^n,\ 1-t=-v^n$. It induces multiplication by $n$
on geometric kernels and retains the arithmetic
translation $(-1,-1)$. Thus the quotient extension
is the arithmetic fundamental group of $T$ even
when $U_n$ has no rational point. This supplies
the comparison underlying (373)--(379).
\end{proof}

\begin{statement}{E.17}{Refuted}{Pairwise separation of cusp sections by the minimum quotient}
For every odd $n$, the rational cusps
$P_0=[0:1:-1]$ and $P_\infty=[1:-1:0]$ admit rational
tangent choices whose sections become conjugate in
the rank-two quotient of (373). Both toric Kummer
classes are $(1,-1)$ in E.16. For the three punctures
of $U_1$ the complete families of toric images are
$(w,1)$, $(1,-w)$ and $(w^{-1},-w^{-1})$, with
$w\in\mathbb Q^*$, in the conventional coordinates
$(t,1-t)$ and the tangent coordinates specified below.
\end{statement}
\begin{proof}
Use $u=x_0/x_\infty,v=x_1/x_\infty$. Near $P_0$
take $s=u$; the branch is
$v=-(1+s^n)^{1/n}$ with constant coefficient $-1$.
The unit tangent has leading torus coefficients
$(1,-1)$. Near $P_\infty$ take $s=x_\infty/x_0$;
then $(u,v)=(s^{-1},-s^{-1}(1+s^n)^{1/n})$,
giving the same leading coefficients. Adjoining
compatible roots of these units computes the
tangential cocycle: coefficientwise Galois action
fixes the formal parameter and acts on roots of
the leading coefficients. Both classes therefore
coincide. This also derives the applicable case of
Wickelgren's tangential formula, Lemma 4
\cite{E_Wickelgren}.

For $U_1$ in coordinate $t$, at $0$ with tangent
$w$ the units $t,1-t$ give $(w,1)$; at $1$ with
coordinate $t-1$ they give $(1,-w)$; at infinity
with coordinate $1/t$ they give
$(w^{-1},-w^{-1})$. These exhaust rational
tangential choices. E.16 determines equality of
their images. The unit tangent at $0$ and tangent
$-1$ at $1$, for example, both give $(1,1)$.
The exhibited coincidences are recorded in
\texttt{e\_obstructions.py}. They concern precisely
this quotient, with the quotient map specified.
\end{proof}

\begin{statement}{E.18}{Proved}{The real obstruction for the manuscript quartic}
The geometrically abelianized arithmetic fundamental
group of $U_4/\mathbb Q$ has no continuous section.
The section-class sets are also empty for its
geometrically metabelian, full profinite and
geometric pro-$2$ quotients. The obstruction
already lies in a group of order two in
$H^2(C_2,H_1(U_4(\mathbb C),\mathbb Z_2))$.
\end{statement}
\begin{proof}
The plus quartic has no real points, by positivity
of fourth powers. Complex conjugation acts freely
on $U_4(\mathbb C)$. Its quotient is a connected
nonorientable punctured surface of Euler
characteristic $-8$, since the double cover has
Euler characteristic $-16$. It retracts to a
graph, so its fundamental group $J$ is free of
rank nine. The geometric subgroup $N$ of index
two is free of rank seventeen. Finite complex
topological covers with conjugation descent
data identify the real arithmetic extension
with the profinite completion of
$1\to N\to J\to C_2\to1$: descend the complex
cover with its antilinear involution; in the
other direction analytify a finite etale cover.
This is the equivariant finite-cover comparison,
and finite-index completion preserves this
extension.

Free-generator changes put $J\to C_2$ in the
form $a\mapsto\bar1$, $b_i\mapsto\bar0$
for $1\leq i\leq8$. A Schreier basis of $N$ is
\[
 a^2,\quad b_i,\quad ab_i a^{-1}
 \quad(1\leq i\leq8).
\]
Conjugation by $a$ fixes the first abelian basis
vector and exchanges each remaining pair. The
factor set at $(\bar1,\bar1)$ is $a^2$.
Replacing the lift by $na$ changes its square
by $(1+\operatorname{Ad}a)[n]$. Its first
coordinate is $1+2n_0$, which cannot vanish
in $\mathbb Z_2$. Each exchanged pair has
invariant diagonal equal to the image of
the norm. Hence the indicated cohomology
group is $\mathbb Z/2$, and the extension
is its nonzero class. The finite program
tests the factor-set equation for every
modulus dividing twenty-four and for
additional odd powers; precisely the
even moduli are obstructed.

A rational section restricts to a real
decomposition group, and each quotient
in the statement projects to this 2-adic
abelian quotient. All are therefore
obstructed. Its image under
$N^{\mathrm{ab}}\to4\widehat{\mathbb Z}(1)^2$
is zero in $H^2$: conjugation acts as
$-1$ on that torsion-free rank-two module,
whose invariants vanish. This explicit
pushout explains the splitting of (373)
together with its failure to lift.
The odd-adic and rational pushouts of
the first-coordinate obstruction also
vanish because two is invertible there;
this cohomology calculation makes no
assertion of a global section there.
\end{proof}

\begin{statement}{E.19}{Proved}{The cubic cusps are separated by their own 3-adic abelianization}
The three rational cusps $[0:1:-1]$,
$[1:0:-1]$, $[1:-1:0]$ determine pairwise
nonconjugate sections of the geometric
pro-$3$ abelianization of
$\pi_1(U_3/\mathbb Q)$, for every rational
tangent choice. They remain distinct in
the geometric abelian, metabelian and
full profinite quotients.
\end{statement}
\begin{proof}
Filling the cusps maps the affine
fundamental group to the Tate module of
the elliptic curve $X_3$, based at one
cusp. A cusp section maps to the point
section there, independently of its
tangent. Each difference of two distinct
cusps has exact order three: the ratio
of their tangent lines has divisor
three times their difference, and a
genus-one curve has no function with a
single simple pole. Section C gives
the same exact-order calculation.

The cubic has good reduction at two.
Its equation over $\mathbb F_2$ has
exactly the three projective points
with two nonzero coordinates, as the
program enumerates. The kernel of
good reduction over $\mathbb Q_2$ is
pro-$2$, from its formal-parameter
congruence filtration. Odd torsion
therefore injects into this group of
three points. A rational point of
exact order three cannot be divisible
by three: a preimage would have order
nine. Its Kummer class modulo three
is nonzero. The difference of the
two point sections is this class
in $H^1(G_{\mathbb Q},X_3[3])$.
Projection from affine 3-adic homology
detects it and proves the assertion.
This map also identifies information
lost in the quotient of E.17.
\end{proof}

\begin{statement}{E.20}{Proved}{A new obstruction to lifting a toric Kummer section}
For $U_1$ based at the tangent from $0$
towards $1$, the abelian section with
Kummer coordinates $(2,3)$ does not
lift to the class-two quotient with
geometric group
\[
 \langle x,y,z\mid x^2=y^2=z^2=1,
 [x,y]=z,\ [x,z]=[y,z]=1\rangle.
\]
Its obstruction is the nonzero class
$\{2\}\cup\{3\}$ in
$H^2(G_{\mathbb Q},\mu_2)$.
\end{statement}
\begin{proof}
The eight normal forms $x^a y^b z^c$
have central multiplication term
$b a'$. The code verifies this group
law and presentation. The tangential
Galois action is trivial on it:
$x\mapsto x^\chi=x$ and
$y\mapsto f^{-1}y^\chi f=y$, since
$\chi$ is odd and the commutator
$f$ is central here. Pulling the
central factor set back along the
two Kummer characters gives their
cup product (the order of factors
does not change its class modulo
two).

Over $\mathbb Q_3$, the quadratic
extension generated by $\sqrt2$ is
unramified because $2$ is a
nonsquare modulo three. Norms have
even valuation, so $3$ is not a
norm. The crossed-product
description of the cup product
identifies its vanishing with
the norm equation
$3=u^2-2v^2$ over $\mathbb Q_3$:
represent $\sqrt2$ by its
multiplication matrix, and the
matrix equation for the second
quaternion generator is exactly
this norm equation. Thus the
local and the global class are
nonzero. A lifted section would
give a cochain whose coboundary
is this factor set, a contradiction.
\end{proof}

\begin{statement}{E.21}{Refuted}{The unmarked depth assignment fails to descend}
The quotient of $\mathbb L(x,y)$
by $D_y^3$, with its canonical
projection, is not functorial for
all automorphisms of the unmarked
free group. It is functorial for
maps preserving the ideal of the
marked puncture and its lower
central filtration. Thus E.5
tests the five specified ordinary
group approximations; the depth
construction additionally uses
this marking and these morphisms.
\end{statement}
\begin{proof}
The swap of $x,y$ sends
$[y,[y,[y,x]]]\in D_y^3$ to
$[x,[x,[x,y]]]$, a nonzero
depth-one element of E.2.
It therefore fails to preserve
the kernel of the canonical
projection, which proves the
non-descent. A map preserving
the marked ideal preserves its
lower central terms and does
induce the quotient map.
No LM8 return specifying an
additional unmarked functor is
present in the supplied archive.
\end{proof}

\begin{statement}{E.22}{Refuted}{The scheme coordinates do not carry the requested relative action}
Put $q_2:F\to(\mathbb Z/2)^2$ and extend the relative completion
of E.14 to $\mathbb Q_\ell$. There is no homomorphism from the
normalized diagonal arithmetic image in
$\mathcal M(0)(\widehat{\mathbb Z})$ to
$\operatorname{Out}(\mathcal R_{q_2,\mathbb Q_\ell})$
whose composite with $G_{\mathbb Q}$ is its actual outer action.
In particular no such genuine action descends from those coordinates.
The actual action is defined on the continuous geometric
fundamental group, and its relative universal property gives its
action on $\mathcal R_{q_2,\mathbb Q_\ell}$. E.15 gives the
faithful discrete version and its exact image.
\end{statement}
\begin{proof}
First we type the comparison with the arithmetic action. The continuous
unipotent completion of the free profinite group $\widehat N_2$
over $\mathbb Q_\ell$ is the free prounipotent group on its five
Schreier generators. Indeed each finite tuple in a unipotent
$\mathbb Q_\ell$-group generates a relatively compact subgroup,
as follows from upper-triangular coordinates and their polynomial
group law. Its homomorphism from the discrete free group consequently
extends continuously to the free profinite group. The two universal
properties identify this completion with
$N_2^{\rm un}\otimes\mathbb Q_\ell$. The transversal construction
in E.14 then gives the continuous relative completion and its
functorial arithmetic action. The normalized Galois action is the
identity on $(\mathbb Z/2)^2$, because its abelian multiplier is
$\chi\equiv1\pmod2$.

The homomorphism $F\to G\times F^{\rm un}$, $g\mapsto(q_2(g),g)$,
and the relative universal property give the canonical surjection
\[
 \mathcal R_{q_2,\mathbb Q_\ell}\longrightarrow
 F^{\rm un}_{\mathbb Q_\ell}.
\]
It is the maximal unipotent quotient, so its kernel is characteristic.
Every genuine automorphism descends through this map and inner
automorphisms descend to inner automorphisms. This gives the needed
map of outer automorphism groups. A nontrivial outer action on its
target is therefore a nontrivial outer action on its source.

A.19 proves that the normalized arithmetic image is contained in
$\mathcal M(0)\simeq\Gamma$. On $G_{\mathbb Q^{\rm cyc}}$ its
multiplier-scheme coordinate is the identity, and the exact law is
that of $C_2\times\widehat{\mathbb Z}$, an abelian group.
Every commutator of two elements of $G_{\mathbb Q^{\rm cyc}}$
thus maps to the identity in these scheme coordinates.

We now exhibit a commutator whose actual outer unipotent action is
nontrivial. E.3 gives nonzero actual graded Galois classes in
weights three and five, with depth-one coefficients proportional
to $\bar\sigma_3,\bar\sigma_5$. Choose representatives $g_3,g_5$
of nonzero such classes and choose $\tau\in G_{\mathbb Q}$ with
$\chi_\ell(\tau)=1+\ell$. The cyclotomic action on weight $w$
is multiplication by $\chi_\ell^w$. Hence
$h_w=[\tau,g_w]$, for $w=3,5$, has leading class
$((1+\ell)^w-1)$ times that of $g_w$, which is nonzero over
$\mathbb Q_\ell$. Both $h_w$ lie in
$G_{\mathbb Q^{\rm ab}}=G_{\mathbb Q^{\rm cyc}}$, since they
are commutators. Their commutator has nonzero weight-eight,
depth-two outer class: its polynomial is a nonzero scalar multiple
of $B_{2,4}$ computed in E.3. This assertion uses actual graded
Galois classes, not arbitrary special derivations.

Consequently $[h_3,h_5]$ is killed by $\mathcal M(0)$ and acts
nontrivially on $\operatorname{Out}(F^{\rm un}_{\mathbb Q_\ell})$,
and thus on the relative completion. A homomorphism agreeing with
both actions would send the identity to this nonidentity element,
which proves the claimed non-descent. For a separate genuine
geometric witness, $w=[x,y]\in\Pi_0'$ has trivial final coordinates
$(a,b,U)=(1,1,\bar w)$ but its inner action on the relative
completion is nontrivial by E.15. The arithmetic witness above
also survives passage to outer automorphisms.
\end{proof}

\begin{statement}{E.23}{Open}{The remaining cubic section calculation}
The complete set of continuous section classes for the geometric
metabelian and marked depth-two arithmetic quotients of
$U_3/\mathbb Q$ is not determined in this return. The remaining
calculation is the image and every fiber of their section maps
to E.16. The first unresolved concrete nonabelian case here is
the class-two pro-$3$ quotient of the cubic's rank-ten geometric
fundamental group, with its integral arithmetic action.
\end{statement}
\noindent\emph{Work and obstruction.}
E.18 determines the empty set for the plus quartic and all
quotients retaining its $2$-adic geometric abelianization;
E.25 extends this to every even plus-Fermat level. E.24 computes
the entire cubic abelian section set and every fiber over the
minimum torus. E.19 separates the three rational cubic cusps
inside that set, while E.17 proves their collision under the
stated toric map.

The relation-module presentation E.1, the rank-two factor-set
calculation, local cup products, the exact generalized-Jacobian
extension, and a new three-punctured-line map on the cubic were
successive attempts to compute the nonabelian fibers. They give
the following actual refinements. E.28 exhibits a whole
nonempty abelian toric fiber with no class-two lifts. E.29
classifies every rational toric pair for the inherited cubic
class-two quotient. E.32 completes that calculation for every
completed pro-$3$ Kummer pair, including the exact convergence
condition and all three classes in each nonempty fiber; it also
exhibits an actual obstruction invisible at every finite level.
E.31 computes a genuine intrinsic quotient with its
arithmetic origin term: in one toric fiber the zero abelian
class lifts, whereas an explicit other class and its entire
congruence coset have no intrinsic class-two lift. Thus the
two minimum coordinates do not determine intrinsic lifting,
as witnessed by this exact same-fiber comparison.
E.34 computes the entire intrinsic abelian family unramified
outside three and proves that every one of its classes lifts
to the full geometric pro-$3$ group.
E.35 evaluates every section of the intrinsic finite
Heisenberg quotient by a cubic norm equation and gives all
of its nonempty fibers in explicit Kummer coordinates.
E.36 proves that every abelian class lifts through the
cuspidally central quotient retaining all ten abelian
coordinates; each fiber is the specified completed group
of rational units. Its quotient map from the full
intrinsic class-two group has a central kernel of rank $44$.

These calculations use the quotient homomorphisms specified in
E.25. They do not supply every obstruction and every fiber for
the full rank-ten intrinsic class-two group in the statement.
After E.36 the next calculation is the actual obstruction
in that rank-$44$ kernel, including its integral lattice
and continuous inverse-limit conditions.
Wickelgren's opened formulas apply directly to the ambient
rank-two group; the additional intrinsic quotient in E.31
requires its separately computed arithmetic action. E.22
settles the proposed descent from the manuscript scheme by an
actual Galois counterexample. The next unresolved step is to
evaluate the full intrinsic central extension on every class
of the explicitly computed set E.24, retaining the origin
terms and the geometric directions killed by E.25's quotient
map. A name for that nonabelian cohomology is not its evaluation.
